% Standalone N=6 CHL, K3 quotient, and Borcherds-denominator data.
% Primary sources: Mukai 1988; Hashimoto 2012; Cheng--Duncan--Harvey
% 2014; Gritsenko--Nikulin 1995, 1998; Gritsenko 1999;
% Clery--Gritsenko 2013.

\providecommand{\kBKM}{\kappa_{\mathrm{BKM}}}
\providecommand{\kch}{\kappa_{\mathrm{ch}}}
\providecommand{\kcat}{\kappa_{\mathrm{cat}}}
\providecommand{\kfib}{\kappa_{\mathrm{fiber}}}
\providecommand{\K}{\mathrm{K3}}
\providecommand{\Fix}{\mathrm{Fix}}
\providecommand{\Z}{\mathbb{Z}}
\providecommand{\Q}{\mathbb{Q}}
\providecommand{\C}{\mathbb{C}}

\section*{$N = 6$ CHL data and the cyclic K3 quotient}

\paragraph{Source statement.}
At level $N = 6$ the relevant input is the $(g_6,h_6)$-twined weak
Jacobi form of weight $0$ and index $1$ on $\Gamma_0(6)$, with
$h_6=g_6^s$ and $s \in \{0,1,2,3\}$. The Fourier constants below are
the K3 twining and Borcherds-product constants in the CHL
normalisation.

\paragraph{Symplectic order-six K3 data.}
Let $S$ be a K3 surface and let
$g_6 \in \mathrm{Aut}_s(S)$ have order $6$. Mukai 1988,
\emph{Invent.\ Math.}~$94$, Table~$2$, gives
\[
  \#\Fix(g_6)=2,\qquad
  \mathrm{rk}\,(H^2(S,\Z)^{g_6})=14,\qquad
  \mathrm{rk}\,(H^2(S,\Z)_{g_6})=8 .
\]
Nikulin's fixed-point counts for the powers are
\[
  \#\Fix(g_6)=2,\qquad
  \#\Fix(g_6^2)=6,\qquad
  \#\Fix(g_6^3)=8 .
\]
The cyclic quotient singularities of $S/\langle g_6\rangle$ follow from
stabiliser orbits. The two points fixed by the full group
$\Z/6$ give two $A_5$ singularities. The remaining four
$\langle g_6^2\rangle \cong \Z/3$ fixed points form two
$\Z/6$-orbits, giving two $A_2$ singularities. The remaining six
$\langle g_6^3\rangle \cong \Z/2$ fixed points form two
$\Z/6$-orbits, giving two $A_1$ singularities. Hence
\[
  S/\langle g_6\rangle
  \quad\hbox{has Du-Val type}\quad
  2A_5+2A_2+2A_1 .
\]
Burnside gives
\[
  e(S/\Z_6)=\frac{24+2+6+8+6+2}{6}=8,
\]
and the exceptional curves contribute
$2\cdot5+2\cdot2+2\cdot1=16$, so the crepant resolution has Euler
characteristic $24$.

\paragraph{Umbral Niemeier anchor.}
Cheng--Duncan--Harvey 2014 assigns the $N = 6$ umbral moonshine case to
the Niemeier lattice
\[
  N_6 = D_4^{\oplus 6}
\]
with root system $6D_4$, Coxeter number $6$, and umbral group
$G^{(6D_4)} = 3.\mathrm{Sym}_6$ of order $2160$. The distinct Niemeier
lattice $A_5^4D_4$ has the same Coxeter number and outer group of order
$144$. The $6A$ CHL/umbral anchor is $6D_4$. The ADE labels in
$2A_5+2A_2+2A_1$ describe local surface singularities on
$S/\langle g_6\rangle$; the label $6D_4$ describes a global even
unimodular rank-$24$ lattice.

\paragraph{Twined Jacobi seed.}
For the Mathieu class $6A$ with cycle shape
$1^2 2^2 3^2 6^2$, the K3-twined weak Jacobi form in the
Eichler--Zagier normalisation satisfies
\[
  \phi^{(g_6)}_{0,1}(\tau,z)
  =\bigl(\zeta^{-1}+2+\zeta\bigr)+q(\cdots),
  \qquad \zeta=e^{2\pi iz}.
\]
Thus the Borcherds-input constant is
\[
  c_6(0)=[q^0\zeta^0]\,\phi^{(g_6)}_{0,1}=2 .
\]
This is the constant entering the BKM denominator weight. The Niemeier
Coxeter number and the quotient Euler characteristic are separate
entries in the level-$6$ data.

\paragraph{Borcherds denominator.}
At the primitive CHL/BKM-denominator scope,
Gritsenko's additive lift and the Borcherds weight theorem give
\[
  \Delta_1^{(6)}
  =\mathrm{BorLift}_6\bigl(\phi^{(g_6)}_{0,1}\bigr),
  \qquad
  \mathrm{wt}\,\Delta_1^{(6)}
  =\kBKM(\Delta_1^{(6)})
  =\frac{c_6(0)}{2}=1 .
\]
Equivalently, on the CHL ladder
\[
  \bigl(c_N(0)\bigr)_{N\in\{1,2,3,4,6\}}=(10,8,6,4,2),
  \qquad
  \bigl(\kBKM(\Phi_N)\bigr)=(5,4,3,2,1).
\]

\paragraph{$\kappa_\bullet$ values at level $6$.}
For the K3 product geometry,
\[
  \kcat(\K\times E)=
  \chi(\mathcal{O}_{\K})\chi(\mathcal{O}_{E})=2\cdot0=0 .
\]
For the level-$6$ denominator,
\[
  \kBKM(\Delta_1^{(6)})=1 .
\]
The fibre-lattice datum attached to the order-six symplectic class is
$\kfib^{(6)}=14$, the invariant lattice rank. The chiral invariant
$\kch$ is determined by the chosen $\Phi$-image and is independent of the
Borcherds weight.

\paragraph{Summary.}
\begin{center}
\footnotesize
\begin{tabular}{l|c|c|c|c|c}
object & $\#\Fix(g_6)$ & invariant rank & $c_6(0)$ & $\kBKM$ & quotient type\\
\midrule
$6A$ symplectic K3 & $2$ & $14$ & $2$ & $1$ & $2A_5+2A_2+2A_1$\\
\end{tabular}
\end{center}

Primary references: Mukai 1988 \emph{Invent.\ Math.}~$94$, Table~$2$;
Hashimoto 2012 \emph{Tohoku Math.\ J.}~$64$; Cheng--Duncan--Harvey 2014
\texttt{arXiv:1307.5793}, Tables~$1$--$2$; Gritsenko--Nikulin 1995
Part~II; Gritsenko--Nikulin 1998 \emph{Amer.\ J.\ Math.}~$119$;
Gritsenko 1999 \emph{St.\ Petersburg Math.\ J.}~$11$; Clery--Gritsenko
2013 \emph{Proc.\ London Math.\ Soc.}~$108$.
